\section{Eredmények számítása és megjelenítése}

A bemeneti adatok feladatonként címkézett, idősoros mérési pontokból állnak,
amelyek tartalmazzák a talpi nyomáscsatornák (\texttt{ch0, ch1, ch2}) nyers értékeit,
a normalizált háromtengelyű gyorsulási adatokat, valamint a rendszer-időbélyeget.


\textbf{\subsection{Statikus feladatok jelfeldolgozása}}

A statikus gyakorlatok (pl. kétlábas állás, tandem állás) során az algoritmus
elsősorban az egyensúlyi helyzet és a stabilitás megtartását vizsgálja
\cite{posturalcontrol}.

\subsubsection{Súlypontvándorlás (Center of Pressure - CoP)}

A súlypont kilengését a talp elülső része (lábujjak és talppárna) és a sarok közötti,
illetve a lábfej jobb és bal oldala közötti nyomáskülönbségek alapján számítom.
A vertikális tengelyen az algoritmus megvizsgálja, hogy a statikus mérés során a
sarokterhelés átlagos aránya meghaladja-e a 65\%-ot a teljes talpnyomáshoz
viszonyítva. Mivel a Parkinson-kóros betegekre gyakran jellemző a súlypont akaratlan
hátrahelyeződése, így a hátraesés kockázata mérhető. A horizontális és vertikális
kilengések  varianciájának összeadásával egy stabilitási mutatót kapunk.

\subsubsection{Remegési index (Tremor)}

A nyugalmi remegés kimutatására az algoritmus a háromdimenziós gyorsulásvektorok
euklideszi normájának (hosszának) varianciáját számítja ki.
A rendszer a statikus gravitációs hatást az átlagos gyorsulás levonásával vonja ki jelből,
így az eredmény a láb mikro-rezgéseinek intenzitását reprezentálja.

\subsection{Dinamikus feladatok jelfeldolgozása}

A mozgást igénylő feladatoknál (például helyben járás, séta) a fókusz a lépésciklusokra
és a mozgáskivitelezés minőségére helyeződik \cite{gaitanalysis}.

\subsubsection{Lépésdetektálás és sarokleütés (Heel Strike)}

A lépéseket a függőleges (z-irányú) gyorsulási adatok lokális maximumainak keresésével
 detektálja az algoritmus, egy küszöbérték (550) alkalmazásával. A szenzorzajból
eredő hamis dupla detektálások elkerülése érdekében a rendszer egy 15 mintavételi
ciklusos holtidőt alkalmaz a lépések között. A lépés talajra érésének pillanatában a
talpi nyomáseloszlás arányaiból kiszámítható a sarok érkezésének módja.
Ez megmutatja, hogy a felhasználó a sarkán gördül-e, vagy a Parkinson-kórra
jellemző csoszogó járást mutatja, ahol az előtalp terhelése a jellemző.

\subsubsection{Fáradási index}

A fáradás méréséhez a rendszer a munkamenet adatsorát megfelezi, majd kiszámítja
az első és a második félidőben mért 3D gyorsulási vektorok átlagos energiáját (hosszát).
A két érték közötti százalékos amplitúdócsökkenés megmutatja, hogy a gyakorlat végére
mennyire esett vissza a végtagmozgás dinamikája.

\subsection{Lefagyás (Freezing of Gait) detektálása}

A Freezing of Gait (FoG) egy hirtelen, átmeneti epizód, amikor a beteg képtelen
a járás folytatására. Detektálására az algoritmus egy 50 mintavételi pontból álló,
folyamatosan csúszó időablakot alkalmaz az adatsoron \cite{fogdetection}.
A lefagyást két egyidejű feltétel teljesülése azonosítja: a horizontális gyorsulások
varianciája egy adott küszöbérték fölé ugrik (magas frekvenciájú lábremegés jelentkezik),
miközben a maximális és minimális nyomásértékek különbsége minimális marad
(a súlypont nem mozdul el, tehát nincs tényleges lépés).

\section{Eredmények vizualizációja}

A feldolgozott eredmények megjelenítése a Jetpack Compose rajzoló API-ja (\texttt{Canvas})
segítségével történik, egyedi, hardveresen gyorsított grafikonokon.

\begin{itemize}
\item Talpi nyomáseloszlás (Hőtérkép): A három nyomásszenzor átlagos
értékeit egy lábfejet formázó alakzaton jeleníti meg, ahol a színátmenetek
(zöldtől a pirosig)  jelzik a terhelés mértékét.
\item Egyensúlyi helyzet (CoP Scatter Plot): Egy céltáblát formázó szórásdiagram,
amely a nyomáskülönbségek alapján pontokként rajzolja ki a feladatok alatti horizontális
és vertikális súlypontvándorlást.
\item \textbf{Lépésciklusok (Ritmus-vonal):} A függőleges gyorsulás utolsó
100 mérési pontját egy vonaldiagram ábrázolja, amely vizuálisan megmutatja a
lépések ritmusát, egyenletességét vagy az esetleges aszimmetriát.
\end{itemize}

\textbf{\subsection{A vizualizáció és az elemzés célja}}

A vizualizáció célja nem egy diagnosztikai eszköz létrehozása volt, hanem inkább példákon
keresztül bemutatni, hogy a gyakorlatok során gyűjtött adatok hogyan értelmezhetők és milyen
információkat hordozhatnak.